\documentclass[12pt]{article}

% === Packages ===
\usepackage[margin=1in]{geometry}
\usepackage{amsmath, amssymb}
\usepackage{setspace}
\usepackage{titlesec}
\usepackage{enumitem}
\usepackage{helvet}
\renewcommand{\familydefault}{\sfdefault} % modern sans-serif look

% === Formatting tweaks ===
\setstretch{1.15}
\setlength{\emergencystretch}{3em}
\setlist[enumerate]{itemsep=0.4em, topsep=0.2em}
\titleformat{\section}{\large\bfseries\sffamily}{\thesection.}{1em}{}
\titleformat{\subsection}{\normalsize\bfseries\sffamily}{\thesubsection}{1em}{}

% === Title ===
\title{\vspace{-2em}\textbf{Applied Mathematics Oral Exam Questions and Model Answers}\\
\large Nondimensionalization and Similarity}
\author{Prepared for Graduate Oral Examination}
\date{}

\begin{document}

\maketitle

\section*{Questions and model answers}

% -------------------------------------------
\subsection*{1. Why do we nondimensionalize the Navier--Stokes equations?}

\textbf{Answer.}
Nondimensionalization allows us to:
\begin{itemize}
    \item compare flows of different physical scales,
    \item reduce the number of free parameters,
    \item identify the dominant physical balances,
    \item improve conditioning and reveal appropriate numerical scales.
\end{itemize}
It does not guarantee numerical stability; stability still depends on the
spatial discretization, time integrator, and coupling algorithm. A reproducible
IBFE model should state either a complete dimensional parameter set or a
consistent set of dimensionless groups and reference scales. Typical variables
are $\mathbf{x}^*=\mathbf{x}/L$, $\mathbf{u}^*=\mathbf{u}/U$, and $t^*=tU/L$.


% -------------------------------------------
\subsection*{2. Derive the nondimensional incompressible Navier--Stokes equations.}

\textbf{Answer.}
Start from dimensional Navier--Stokes:
\[
\rho\left(
\frac{\partial \mathbf{u}}{\partial t} + (\mathbf{u}\cdot\nabla)\mathbf{u}
\right)
= -\nabla p + \mu\nabla^2\mathbf{u}.
\]

Choose a characteristic length $L$, speed $U$, advective time $L/U$, and
dynamic-pressure scale $\rho U^2$:
\[
\mathbf{x}^*=\frac{\mathbf{x}}{L},\quad t^*=\frac{tU}{L},\quad \mathbf{u}^*=\frac{\mathbf{u}}{U},
\quad p^*=\frac{p}{\rho U^2},\quad \nabla^*=L\nabla.
\]
Substituting and rearranging yields
\[
\frac{\partial \mathbf{u}^*}{\partial t^*}
+ \mathbf{u}^*\cdot\nabla^*\mathbf{u}^*
= -\nabla^* p^* + \frac{1}{Re}\nabla^{*2}\mathbf{u}^*,
\qquad \nabla^*\cdot\mathbf{u}^*=0,
\]
where
\[
Re = \frac{\rho UL}{\mu}=\frac{UL}{\nu}.
\]
The pressure scale is a choice: a viscous scale $\mu U/L$ is often more
revealing for creeping flow and produces an equivalent rescaled equation.


% -------------------------------------------
\subsection*{3. What is the physical meaning of the Reynolds number in IBFE simulations?}

\textbf{Answer.}
For stated choices of $U$ and $L$, $Re=\rho UL/\mu$ compares inertial and
viscous scales. It is a control parameter, not a complete label for the
solution. In IBFE:
\begin{itemize}
    \item Low $Re$ favors a viscous or creeping-flow balance; explicit treatment
          of diffusion can impose a time-step restriction, whereas an implicit
          fluid solve changes that restriction.
    \item High $Re$ can produce thin shear layers and wakes that require adequate
          spatial and temporal resolution; geometry and disturbances still
          determine whether transition occurs.
    \item Structural response is governed by additional dimensionless stiffness,
          mass, and damping parameters. Changing $Re$ alone does not determine
          how strong the immersed force should be.
\end{itemize}
Biological and engineered systems span many orders of magnitude in $Re$, so the
reference scales and the regime of interest must always be reported.
% -------------------------------------------
\subsection*{4. How do you nondimensionalize force per unit length in IBFE?}

\textbf{Answer.}
First specify the measure associated with the force. Suppose $\mathcal{F}$ is a
two-dimensional force resultant per unit out-of-plane depth, with units
$\mathrm{N\,m^{-1}}$. Two consistent normalizations are
\[
\mathcal{F}^{(I)}=\frac{\mathcal{F}}{\rho U^2L},
\qquad
\mathcal{F}^{(V)}=\frac{\mathcal{F}}{\mu U},
\qquad
\mathcal{F}^{(V)}=Re\,\mathcal{F}^{(I)}.
\]
There is no discontinuous switch at $Re=1$: choose the inertial or viscous scale
that exposes the desired balance and use it consistently. A local line-force
density, a total force, and a three-dimensional surface traction have different
units and therefore different scale factors.

% -------------------------------------------
\subsection*{5. How is the dimensionless beam bending stiffness $k_b^{(I)}$ derived?}

\textbf{Answer.}
Under the tutorial's two-dimensional, per-unit-depth beam convention, the
coefficient $k_b$ has units $\mathrm{N\,m}$. Its inertial
normalization is
\[
k_b^{(I)}=\frac{k_b}{\rho U^2 L^3}.
\]
For example, a fourth-derivative bending law has a local force-density scale
$k_b/L^3$, which is compared with the inertial stress scale
$\rho U^2$. With viscous normalization,
\[
k_b^{(V)}=\frac{k_b}{\mu U L^2}.
\]
These dimensionless coefficients differ by a factor of $Re$ but describe the
same dimensional model. A three-dimensional flexural rigidity $EI$ has units
$\mathrm{N\,m^2}$; the convention must not be inferred from the word ``beam''
alone.

% -------------------------------------------
\subsection*{6. Estimate bending stiffness and explain how it affects numerical stability.}

\textbf{Answer.}
The value $k_b^{(I)}\approx 10$ in the tutorial is a problem-specific
starting point, not a universal range for ``moderate'' bending. In IBFE:
\begin{itemize}
    \item A large coefficient relative to fluid and structural scales can make
          explicit coupling stiff, requiring a smaller time step or implicit
          treatment.
    \item A small coefficient permits greater deformation; whether that is
          unrealistic depends on the material, loads, geometry, and boundary
          conditions rather than on a numerical threshold.
    \item Estimate the dimensional rigidity from material/section data or a
          measured force--deflection response, then convert it using the same
          $U$ and $L$ used for the fluid equations.
\end{itemize}

% -------------------------------------------
\subsection*{7. How do you estimate a spring constant for immersed boundary structures?}

\textbf{Answer.}
For a spring-like constraint, an order-of-magnitude effective stiffness is
\[
K_{\text{eff}} \sim \frac{F_{\text{char}}}{\delta_{\max}},
\]
where $\delta_{\max}$ is the acceptable displacement. At moderate or high
$Re$, estimate a physical inertial load using
\[
F_{\text{char}}\sim \tfrac12 C_D\rho U^2A,
\]
using a physical projected area $A$ (or projected length per unit depth in 2D),
not a grid-cell area. At low $Re$, use the appropriate geometry-dependent
viscous drag law. Finally convert $K_{\text{eff}}$ using the units of the actual
node, line, surface, or volume force law implemented in the code.

% -------------------------------------------
\subsection*{8. What dimensionless groups govern oscillatory or pumping flows?}

\textbf{Answer.}
Several related groups appear in the tutorial:
\begin{itemize}
    \item \textbf{Dimensionless time:} $t^*=tU/L$.
    \item \textbf{Dimensionless frequency:} $f^*=fL/U$, the ratio of the
          advective time $L/U$ to the oscillation period $1/f$.
    \item \textbf{Strouhal number:} commonly $St=fL/U$; propulsion studies often
          use the amplitude-based convention $St_A=fA/U$. The often-cited
          $0.2\lesssim St_A\lesssim0.4$ band is empirical and context-dependent,
          not a universal efficiency criterion.
    \item \textbf{Womersley number:} $\alpha=L\sqrt{\omega/\nu}$, where for a
          tube the length scale is the radius, $L=R$; it compares unsteady
          inertia with viscous diffusion in oscillatory flow. Other geometries
          require a clearly stated length.
\end{itemize}
These determine how oscillations couple to fluid inertia and viscosity, crucial
in pumping or swimming IBFE simulations.

% -------------------------------------------
\subsection*{9. How does nondimensionalization guide mesh and time-step choices?}

\textbf{Answer.}
The dimensionless groups determine:
\begin{itemize}
    \item spatial scales that must be resolved (for example, thin shear or
          oscillatory boundary layers at large $Re$ or $\alpha$),
    \item advective, diffusive, elastic, and wave time scales relevant to the
          stability limits of the chosen time integrator,
    \item stiffness in structural models for large $k_b^{(I)}$ or
          $k_s^{(I)}$,
    \item time-step requirements for coupled fluid–structure interaction.
\end{itemize}
For fixed $U$ and $L$, a large dimensionless forcing frequency demands temporal
resolution of the imposed oscillation. Low $Re$ does not inherently require a
fine grid; geometry and solution gradients set spatial resolution, while an
explicit discretization of strong diffusion may restrict the time step.

% -------------------------------------------
\newpage
\subsection*{10. Explain the physical meaning of the Péclet number and how it applies in IBFE.}

\textbf{Answer.}
\[
Pe = \frac{\text{advective transport rate}}
         {\text{diffusive transport rate}}
       = \frac{LU}{D}.
\]
In IBFE:
\begin{itemize}
    \item For passive scalars or mass/chemical transport, $Pe$ determines
          whether advection or diffusion dominates.
    \item If $Pe \gg 1$, steep gradients require enhanced mesh resolution or
          stabilization.
    \item If $Pe \ll 1$, diffusion smooths fields, potentially relaxing mesh
          requirements.
\end{itemize}
This governs both the physical scalar boundary-layer thickness and the choice
of stable, accurate advection--diffusion discretization. It is distinct from
$Re$; for a common $U$ and $L$, $Pe=Re\,Sc$ for mass transfer and
$Pe=Re\,Pr$ for heat transfer.

\end{document}
